\documentclass[sigconf]{acmart}

\usepackage{xeCJK}
\usepackage{subfigure}
\usepackage{graphicx}
\usepackage{array}
\usepackage{enumitem}
\usepackage{multicol}
\usepackage{algorithm,algorithmic}
\usepackage{color}  
%\definecolor{shadecolor}{named}{Gray}  
\definecolor{shadecolor}{rgb}{0.92,0.92,0.92}  
\usepackage{framed}
\usepackage{ulem}
\usepackage{tikz}
\usetikzlibrary{decorations.markings}
\usepackage{makecell}


%% Font
\CJKfontspec{Noto Serif CJK TC} %思源宋體

%% ----------
\newcommand{\answer}[2]{%
  \if\relax\detokenize{#2}\relax
    \underline{\mbox{請作答}}%
  \else
    \textcolor{blue}{#2}%
  \fi
}


\begin{document}

\title{114學年度第二學期~離散數學HW10}

\author{班級：\underline{資工一}、學號：\underline{114590XXX}、姓名：\underline{蘇東坡}}
\orcid{}
\affiliation{%
  \institution{}
  \city{}
  \country{}
}


%% 刪除ACM Reference Format信息
\settopmatter{printacmref=false} % Removes citation information below abstract
\renewcommand\footnotetextcopyrightpermission[1]{} % removes footnote with conference information in first column
\pagestyle{plain} % removes running headers

\maketitle


%% ----- 作業繳交說明 -----
\section{題目}
\subsection{作業繳交說明}
\begin{shaded}
\begin{itemize}
    \item[-] 從 overleaf 下載 hw10.zip
    \item[-] 從 overleaf 下載 hw10.pdf
    \item[-] 建立一個新檔案夾，命名如下
    \item[-] \color{red}\begin{verbatim}離散數學_114590xxx_姓名_HW10\end{verbatim}
    \item[-] \color{black}放入 hw10.zip
    \item[-] 放入 hw10.pdf
    \item[-] 將此檔案夾，壓縮成 zip 檔,檔名如下
    \item[-] \color{red}\begin{verbatim}離散數學_114590xxx_姓名_HW10.zip\end{verbatim}
    \item[-] \color{black}將此壓縮檔，上傳到北科 i 學園 PLUS
    \item[-] Do not share your homework with any living creatures.
    \item[-] 作業遲交，一天內成績打五折，超過一天以零分計。
    \item[-] 畫圖請使用latex指令，勿手寫。
\end{itemize}
\end{shaded}

%% ----- Question -----
\subsection{題號}
\begin{shaded}
    \begin{itemize}
	\item[-] page 683, chapter 10.1 Exercise 2
	\item[-] page 699, chapter 10.2 Exercise 8
	\item[-] page 710, chapter 10.3 Exercise 8
	\item[-] page 724, chapter 10.4 Exercise 2
	\item[-] page 741, chapter 10.5 Exercise 54
	\item[-] page 753, chapter 10.6 Exercise 26
	\item[-] page 761, chapter 10.7 Exercise 14
	\item[-] page 768, chapter 10.8 Exercise 6
\end{itemize}
\end{shaded}


%% ----- Problem -----
\section{作答}
\subsection{page 683, chapter 10.1 Exercise 2}
\begin{shaded}
    What kind of graph (from Table 1) can be used to model a highway system between major cities where
    \begin{enumerate}[label=(\alph*)]
        \item there is an edge between the vertices representing cities if there is an interstate highway between them?
        \item there is an edge between the vertices representing cities for each interstate highway between them?
        \item there is an edge between the vertices representing cities for each interstate highway between them, and there is a loop at the vertex representing a city if there is an interstate highway that circles this city?
    \end{enumerate}
\end{shaded}
\begin{table}[h]
    \centering
	\caption{Graph Terminology.}
    \begin{tabular}{|p{2.8cm}|p{3.1cm}|p{1.8cm}|p{1.1cm}|}
        \hline
        $\textsl{Type}$ & $\textsl{Edges}$ & $\textsl{\makecell{Multiple Edges \\ Allowed?}}$ & $\textsl{\makecell{Loops \\Allowed?}}$ \\ 
        \hline
        Simple graph & Undirected & No & No \\ 
        Multigraph & Undirected & Yes & No \\ 
        Pseudograph & Undirected & Yes & Yes \\ 
        Simple directed graph & Directed & No & No \\ 
        Directed multigraph & Directed & Yes & Yes \\ 
        Mixed graph & Directed and undirected & Yes & Yes \\ 
        \hline
    \end{tabular}
\end{table}
\begin{enumerate}[label=(\alph*)]
    \item \answer{10.1.2a}{simple graph}. Since there are no parallel edges or loops, and the edges are undirected. 
    \item \answer{10.1.2b}{}
    \item \answer{10.1.2c}{}
\end{enumerate}


\subsection{page 699, chapter 10.2 Exercise 8}
\begin{shaded}
    Determine the number of vertices and edges and find the in-degree and out-degree of each vertex for the given directed multigraph.

    %%2.2圖片
    \begin{center}
        \tikzstyle{point}=[circle,draw=black,fill=black,thick,inner sep=0pt,minimum size=2mm]
        \tikzset{middlearrow/.style={
            decoration={markings,
            mark= at position 0.5 with {\arrow{#1}} ,
            },
            postaction={decorate}
            }
        }
        \begin{tikzpicture}
        \node at ( -1.2,1.2) [point,draw,label=above:\Large$a$] (a) {};
        \node at ( 1.2,1.2) [point,draw,label={[xshift=-8pt]above:\Large$b$}] (b) {};
        \node at ( 1.2,-1.2) [point,draw,label=below:\Large$c$] (c) {};
        \node at ( -1.2,-1.2) [point,draw,label=below:\Large$d$] (d) {};
        \draw[middlearrow={latex},very thick,bend right=20] (b) to (a);
        \draw[middlearrow={latex},very thick,bend right=20] (a) to (b);
        \draw[middlearrow={latex},very thick] (a) to (b);
        \draw[middlearrow={latex},very thick,bend left=20] (b) to (c);
        \draw[middlearrow={latex},very thick,bend right=20] (b) to (c);
        \draw[middlearrow={latex},very thick] (c) to (a);
        \draw[->,>= latex,very thick]  (b) edge [out=-20,in=80,distance=8mm]   (b);
        \draw[->,>= latex,very thick]  (d) edge [out=-20,in=80,distance=8mm]   (d);
        
        \end{tikzpicture}
    \end{center}
    
\end{shaded}

\vspace{1cm}

In this directed multigraph there are \answer{10.2.8a}{} vertices and \answer{10.2.8b}{} edges.\\ The degrees are:\\ 
$\text{deg}^{-}(a) = 2$, $\text{deg}^{+}(a) = 2$, \\
$\text{deg}^{-}(b) = \answer{10.2.8c}{}$, $\text{deg}^{+}(b) = \answer{10.2.8d}{}$, \\
$\text{deg}^{-}(c) = \answer{10.2.8e}{}$, $\text{deg}^{+}(c) = \answer{10.2.8f}{}$, \\
$\text{deg}^{-}(d) = \answer{10.2.8g}{}$, $\text{deg}^{+}(d) = \answer{10.2.8h}{}$.

\clearpage



\subsection{page 710, chapter 10.3 Exercise 8}
\begin{shaded}
    Represent the graph in Exercise 4 with an adjacency matrix.

    %%圖片2.3
    \begin{center}
        \tikzstyle{point}=[circle,draw=black,fill=black,thick,
        inner sep=0pt,minimum size=2mm]
        \tikzset{middlearrow/.style={
                decoration={markings,
                    mark= at position 0.5 with {\arrow{#1}} ,
                },
                postaction={decorate}
            }
        }
        \begin{tikzpicture}
        \node at ( -2,1) [point,draw,label=above:\LARGE$a$] (a) {};
        \node at ( 0,1) [point,draw,label=above:\LARGE$b$] (b) {};
        \node at ( 2,1) [point,draw,label={[xshift=-6pt]above:\LARGE$c$}] (c) {};
        \node at ( -2,-1) [point,draw,label=below:\LARGE$d$] (d) {};
        \node at ( 0,-1) [point,draw,label={[xshift=-9pt]below:\LARGE$e$}] (e) {};
        
        \draw[middlearrow={latex},very thick,bend right=20] (a) to (b);
        \draw[middlearrow={latex},very thick,bend right=20] (b) to (a);
        \draw[middlearrow={latex},very thick,bend right=20] (b) to (c);
        \draw[middlearrow={latex},very thick,bend right=20] (c) to (b);
        \draw[middlearrow={latex},very thick,bend left=20] (a) to (d);
        \draw[middlearrow={latex},very thick,bend left=20] (d) to (a);
        \draw[middlearrow={latex},very thick] (b) to (d);
        \draw[middlearrow={latex},very thick] (b) to (e);
        \draw[middlearrow={latex},very thick] (d) to (e);
        \draw[middlearrow={latex},very thick] (e) to (c);
        \draw[->,>= latex,very thick]  (c) edge [out=-20,in=80,distance=8mm]   (c);
        \draw[->,>= latex,very thick]  (e) edge [out=240,in=340,distance=8mm]   (e);
        \end{tikzpicture}
    \end{center}
    
\end{shaded}



$\begin{bmatrix}
	   0 & 1 & 0 & 1 & 0 \\
       \answer{10.3.8a}{} & \answer{10.3.8b}{} & \answer{10.3.8c}{} & \answer{10.3.8d}{} & \answer{10.3.8e}{} \\
       \answer{10.3.8f}{} & \answer{10.3.8g}{} & \answer{10.3.8h}{} & \answer{10.3.8i}{} & \answer{10.3.8j}{} \\
       \answer{10.3.8k}{} & \answer{10.3.8l}{} & \answer{10.3.8m}{} & \answer{10.3.8n}{} & \answer{10.3.8o}{} \\
       \answer{10.3.8p}{} & \answer{10.3.8q}{} & \answer{10.3.8r}{} & \answer{10.3.8s}{} & \answer{10.3.8t}{} \\
\end{bmatrix}$

\subsection{page 724, chapter 10.4 Exercise 2}
\begin{shaded}
    Does each of these lists of vertices form a path in the following graph? Which paths are simple? Which are circuits? What are the lengths of those that are paths
    \begin{enumerate}[label=(\alph*)]
        \item $a, b, e, c, b$
        \item $a, d, a, d, a$
        \item $a, d, b, e, a$
        \item $a, b, e, c, b, d, a$
    \end{enumerate}

    \begin{center}
        \tikzstyle{point}=[circle,draw=black,fill=black,thick,
        inner sep=0pt,minimum size=2mm]
        \tikzset{middlearrow/.style={
                decoration={markings,
                    mark= at position 0.5 with {\arrow{#1}} ,
                },
                postaction={decorate}
            }
        }
        \begin{tikzpicture}
        \node at ( -2,1) [point,draw,label=above:\LARGE$a$] (a) {};
        \node at ( 0,1) [point,draw,label=above:\LARGE$b$] (b) {};
        \node at ( 2,1) [point,draw,label={[xshift=-6pt]above:\LARGE$c$}] (c) {};
        \node at ( -2,-1) [point,draw,label=below:\LARGE$d$] (d) {};
        \node at ( 0,-1) [point,draw,label={[xshift=-9pt]below:\LARGE$e$}] (e) {};
        
        \draw[middlearrow={latex},very thick,bend left=20] (a) to (b);
        \draw[middlearrow={latex},very thick,bend left=20] (b) to (a);
        \draw[middlearrow={latex},very thick,bend left=20] (a) to (d);
        \draw[middlearrow={latex},very thick,bend left=20] (d) to (a);
        \draw[middlearrow={latex},very thick] (c) to (b);
        \draw[middlearrow={latex},very thick] (b) to (e);
        \draw[middlearrow={latex},very thick] (e) to (d);
        \draw[middlearrow={latex},very thick] (e) to (c);
        \end{tikzpicture}
    \end{center}
\end{shaded}
\begin{enumerate}[label=(\alph*)]
    \item This \uline{~~is~~} (is, is not) a path of length \uline{~~4~~} (if not a path fill 0), it \uline{~~is not~~} (is, is not) a circuit, it \uline{~~is~~} (is, is not) simple.
	
	\item This \answer{10.4.2a}{} (is, is not) a path of length \answer{10.4.2b}{} (if not a path fill 0), it \answer{10.4.2c}{} (is, is not) a circuit, it \answer{10.4.2d}{} (is, is not) simple.

	\item This \answer{10.4.2e}{} (is, is not) a path of length \answer{10.4.2f}{} (if not a path fill 0), it \answer{10.4.2g}{} (is, is not) a circuit, it \answer{10.4.2h}{} (is, is not) simple.

	\item This \answer{10.4.2i}{} (is, is not) a path of length \answer{10.4.2j}{} (if not a path fill 0), it \answer{10.4.2k}{} (is, is not) a circuit, it \answer{10.4.2l}{} (is, is not) simple.
\end{enumerate}

\subsection{page 741, chapter 10.5 Exercise 54}
\begin{shaded}
    A diagnostic message can be sent out over a computer network to perform tests over all links and in all devices. What sort of paths should be used to test all links? To test all devices?
\end{shaded}
\answer{10.5.54a}{} (An Euler, A Hamilton) path will cover every link, so it can be used to test the links.\\
\answer{10.5.54b}{} (An Euler, A Hamilton) path will cover all the devices, so it can be used to test the devices.

\subsection{page 753, chapter 10.6 Exercise 26}
\begin{shaded}
    Solve the traveling salesperson problem for this graph by finding the total weight of all Hamilton circuits and determining a circuit with minimum total weight.

    %%圖片2.6
    \begin{center}
        \tikzstyle{point}=[circle,draw=black,fill=black,thick,inner sep=0pt,minimum size=2mm]
        \begin{tikzpicture}
        \node at ( -1.2,1.6) [point,draw,label=above:\LARGE$a$] (a) {};
        \node at ( 1.2,1.6) [point,draw,label=above:\LARGE$b$] (b) {};
        \node at ( 1.8,-0.4) [point,draw,label=right:\LARGE$c$] (c) {};
        \node at ( 0,-1.6) [point,draw,label=below:\LARGE$d$] (d) {};
        \node at ( -1.8,-0.4) [point,draw,label=left:\LARGE$e$] (e) {};
        \path
            (a) edge[very thick] node [above] {\LARGE 3} (b)
            (b) edge[very thick] node [above,xshift=5pt] {\LARGE 10} (c)
            (c) edge[very thick] node [below,xshift=3pt] {\LARGE 6} (d)
            (d) edge[very thick] node [below,xshift=-3pt] {\LARGE 1} (e)
            (a) edge[very thick] node [above,xshift=-5pt] {\LARGE 7} (e)
            (b) edge[very thick] node [above,xshift=-3pt,yshift=-3pt] {\LARGE 2} (e)
            (e) edge[very thick] node [above,yshift=-1pt] {\LARGE 5} (c)
            (c) edge[very thick] node [above,xshift=3pt,yshift=-3pt] {\LARGE 8} (a)
            (a) edge[very thick] node [above,xshift=3pt,yshift=-4pt] {\LARGE 4} (d)
            (d) edge[very thick] node [above,xshift=-3pt,yshift=-4pt] {\LARGE 9} (b);
        
        \end{tikzpicture}
    \end{center}
    
\end{shaded}

\begin{enumerate}[label=(\alph*)]
    \item Circuit: a-b-c-d-e-a, Weight = 3 + 10 + 6 + 1 + 7 = 27
	\item Circuit: a-b-c-e-d-a, Weight = \answer{10.6.26a}{}
	\item Circuit: a-b-d-c-e-a, Weight = \answer{10.6.26b}{}
	\item Circuit: a-b-d-e-c-a, Weight = \answer{10.6.26c}{}
	
	\item Circuit: a-b-e-c-d-a, Weight = \answer{10.6.26d}{}
	\item Circuit: a-b-e-d-c-a, Weight = \answer{10.6.26e}{}
	\item Circuit: a-c-b-d-e-a, Weight = \answer{10.6.26f}{}
	\item Circuit: a-c-b-e-d-a, Weight = \answer{10.6.26g}{}
	
	\item Circuit: a-c-d-b-e-a, Weight = \answer{10.6.26h}{}
	\item Circuit: a-c-e-b-d-a, Weight = \answer{10.6.26i}{}
	\item Circuit: a-d-b-c-e-a, Weight = \answer{10.6.26j}{}
	\item Circuit: a-d-c-b-e-a, Weight = \answer{10.6.26k}{}
\end{enumerate}
The circuits \answer{10.6.26l}{(start and end at $a$)} and \answer{10.6.26m}{(start and end at $a$)} are the ones with minimum total weight.

\subsection{page 761, chapter 10.7 Exercise 14}
\begin{shaded}
    Suppose that a connected planar graph has 30 edges. If a planar representation of this graph divides the plane into 20 regions, how many vertices does this graph have?
\end{shaded}
Euler's formula says that $v - e + r = 2$.\\
We are given $e = \answer{10.7.14a}{}$ and $r = \answer{10.7.14b}{}$. \\
Therefore $v = 2 - r + e = \answer{10.7.14c}{}$.


\subsection{page 768, chapter 10.8 Exercise 6}
\begin{shaded}
    Find the chromatic number of the given graph.
    
    %%圖片2.8
    \begin{center}
        \tikzstyle{point}=[circle,draw=black,fill=black,thick,
        inner sep=0pt,minimum size=2mm]
        \begin{tikzpicture}
        \node at ( -2,0) [point,draw,label=left:\LARGE$a$] (a) {};
        \node at ( -1,1) [point,draw,label=above:\LARGE$b$] (b) {};
        \node at ( 1,1) [point,draw,label=above:\LARGE$c$] (c) {};
        \node at ( 2,0) [point,draw,label=right:\LARGE$d$] (d) {};
        \node at ( 1,-1) [point,draw,label=below:\LARGE$e$] (e) {};
        \node at ( -1,-1) [point,draw,label=below:\LARGE$f$] (f) {};
        \node at ( 0,0) [point,draw,label=above:\LARGE$g$] (g) {};
        \path
            (a) edge [very thick]  (b)
            (b) edge [very thick]  (c)
            (c) edge [very thick]  (d)
            (d) edge [very thick]  (e)
            (f) edge [very thick]  (a)
            (a) edge [very thick]  (g)
            (b) edge [very thick]  (g)
            (f) edge [very thick]  (g)
            (c) edge [very thick]  (g)
            (d) edge [very thick]  (g)
            (e) edge [very thick]  (g);
        \end{tikzpicture}
    \end{center}
    
\end{shaded}

At least \answer{10.8.6a}{} colors are needed, since ...

\vspace{15cm}

\end{document}
\endinput
%%
%% End of file `sample-sigconf.tex'.